\section*{Capítulo \ref{ch:g1:watching-the-weather} --- Observar el tiempo}

\begin{solution}{exo:g1:watching-the-weather:1}
Soleado: \omwx{sun}\quad lluvioso: \omwx{rain}\quad con nieve: \omwx{snow}\quad
nublado: \omwx{cloud}. Vale cualquier dibujo claro de un sol, una nube con
gotas, una nube con copos de nieve y una nube sola.
\end{solution}

\begin{solution}{exo:g1:watching-the-weather:2}
El parque: sol y calor. La calle: nieve y frío. La colina: viento (y
algunas nubes). La puerta del colegio: niebla.
\end{solution}

\begin{solution}{exo:g1:watching-the-weather:3}
Martes por la mañana: nublado. Viernes por la tarde: chubascos. Domingo por
la tarde: lluvioso.
\end{solution}

\begin{solution}{exo:g1:watching-the-weather:4}
El miércoles: el sol está dibujado por la mañana y por la tarde.
(El jueves por la tarde hay sol y nubes.)
\end{solution}

\begin{solution}{exo:g1:watching-the-weather:5}
La gorra para el sol va con el tiempo de calor y sol; las botas y el
impermeable, con el tiempo lluvioso; la bufanda y los guantes, con el tiempo
frío y con nieve.
\end{solution}

\begin{solution}{exo:g1:watching-the-weather:6}
Hace viento: el viento empuja la bandera.
\end{solution}

\begin{solution}{exo:g1:watching-the-weather:7}
\begin{enumerate}
  \item Falso: el \omterm{def:g1:watching-the-weather:weather}{tiempo} puede cambiar a lo largo del día; por ejemplo, una
        mañana gris y una tarde soleada.
  \item Verdadero.
  \item Verdadero: el Sol brilla tanto que puede dañar el ojo para siempre,
        incluso en un momento.
\end{enumerate}
\end{solution}

\begin{solution}{exo:g1:watching-the-weather:8}
Tardes con lluvia o chubascos: el lunes (lluvia), el viernes (chubascos) y el domingo
(lluvia), o sea, 3. Días fríos: el lunes, el martes y el domingo, o sea, 3.
Los otros 4 días hizo calor, así que la semana fue sobre todo cálida.
\end{solution}

\begin{solution}{exo:g1:watching-the-weather:9}
Mañana: \omwx{cloud} (cielo gris); mediodía: \omwx{rain} (lluvia);
noche: \omwx{suncloud} (cielo azul con unas pocas nubes).
\end{solution}

\begin{solution}{exo:g1:watching-the-weather:10}
Una bandera señala hacia donde va el viento. Si señala hacia la izquierda,
el viento va hacia la izquierda, así que viene de la derecha.
\end{solution}

\begin{solution}{exo:g1:watching-the-weather:11}
Sin hojas: en invierno. Con flores: en primavera. El invierno tiene el
símbolo de la nieve: en invierno hace frío, y en muchos lugares solo nieva
en esa estación.
\end{solution}
